\documentclass[11pt,a4paper]{article}
\usepackage[utf8]{inputenc}
\usepackage{amsmath,amssymb,amsthm}
\usepackage{geometry}
\geometry{margin=1in}

\newtheorem{theorem}{Theorem}
\newtheorem{proposition}{Proposition}
\newtheorem{definition}{Definition}

\title{Euler--Bernoulli--Zeta Relations: Exact Computation}
\author{Valhalla Research Team}
\date{January 2026}

\begin{document}
\maketitle

\begin{abstract}
We verify classical relationships between Euler's Basel problem, Bernoulli numbers, and the Riemann zeta function using exact rational arithmetic (S-wind EPR). All computations are performed with infinite precision, demonstrating exact algebraic identities.
\end{abstract}

\section{The Basel Problem}

\begin{theorem}[Euler, 1735]
\begin{equation}
\zeta(2) = \sum_{n=1}^{\infty} \frac{1}{n^2} = \frac{\pi^2}{6}
\end{equation}
\end{theorem}

\section{Exact Partial Sums}

Using S-wind EPR arithmetic:

\begin{proposition}[Partial Sum]
\begin{equation}
S_N = \sum_{n=1}^{N} \frac{1}{n^2}
\end{equation}
For $N = 20$:
\begin{equation}
S_{20} = \frac{17299975731542641}{10838475198270720}
\end{equation}
\end{proposition}

This is an \textbf{exact fraction}, not a floating-point approximation.

\section{Bernoulli Number Connection}

\begin{theorem}[Euler's Formula for Even Zeta Values]
\begin{equation}
\zeta(2n) = (-1)^{n+1} \frac{(2\pi)^{2n} B_{2n}}{2(2n)!}
\end{equation}
where $B_{2n}$ are Bernoulli numbers.
\end{theorem}

\subsection{First Few Values}

\begin{center}
\begin{tabular}{|c|c|c|}
\hline
$n$ & $\zeta(2n)$ & Exact Form \\
\hline
1 & $\zeta(2)$ & $\pi^2/6$ \\
2 & $\zeta(4)$ & $\pi^4/90$ \\
3 & $\zeta(6)$ & $\pi^6/945$ \\
4 & $\zeta(8)$ & $\pi^8/9450$ \\
\hline
\end{tabular}
\end{center}

\section{Bernoulli Numbers (Exact)}

Using S-wind EPR:

\begin{center}
\begin{tabular}{|c|c|}
\hline
$n$ & $B_n$ \\
\hline
0 & $1$ \\
1 & $-1/2$ \\
2 & $1/6$ \\
4 & $-1/30$ \\
6 & $1/42$ \\
8 & $-1/30$ \\
10 & $5/66$ \\
12 & $-691/2730$ \\
\hline
\end{tabular}
\end{center}

All values are exact fractions.

\section{Verification}

\begin{proposition}
For $n = 1$ (Basel problem):
\begin{equation}
\zeta(2) = \frac{(2\pi)^2 \cdot (1/6)}{2 \cdot 2!} = \frac{4\pi^2 / 6}{4} = \frac{\pi^2}{6}
\end{equation}
\end{proposition}

This confirms Euler's formula using exact algebraic manipulation.

\section{Technical Notes}

\begin{itemize}
    \item All computations use S-wind EPR (infinite precision rationals)
    \item Bernoulli numbers computed via standard recurrence
    \item Partial sums are exact; convergence to $\pi^2/6$ follows analysis
\end{itemize}

\begin{thebibliography}{9}
\bibitem{euler1735} L. Euler, \textit{De summis serierum reciprocarum}, 1735.
\bibitem{bernoulli1713} J. Bernoulli, \textit{Ars Conjectandi}, 1713.
\bibitem{titchmarsh1986} E.C. Titchmarsh, \textit{The Theory of the Riemann Zeta-Function}, Oxford, 1986.
\end{thebibliography}

\end{document}
